% ===== PRÉAMBULE CANONIQUE — à coller VERBATIM en tête de CHAQUE .tex =====
\documentclass[11pt,a4paper]{article}
\usepackage[utf8]{inputenc}\usepackage[T1]{fontenc}\usepackage{lmodern}
\usepackage[french]{babel}
\usepackage{amsmath,amssymb,mathtools}\usepackage{icomma}
\usepackage[version=4]{mhchem}          % \ce{...} : équations & formules
\usepackage{siunitx}\sisetup{locale=FR} % \qty{}{}, \si{}, \ang{}
\usepackage{chemfig}                    % \chemfig{...} : structures développées
\usepackage{xcolor}\usepackage{colortbl}
\usepackage{tikz}\usepackage{pgfplots}\pgfplotsset{compat=1.18}
\usepackage[most]{tcolorbox}\usepackage{enumitem}\usepackage{array,booktabs}
\usepackage[a4paper,margin=2.2cm]{geometry}
\usetikzlibrary{arrows.meta,calc,angles,quotes,positioning,patterns,backgrounds,shapes.geometric}
\definecolor{primary}{RGB}{222,49,99}\definecolor{primarydark}{RGB}{150,22,55}
\definecolor{defcol}{RGB}{0,114,178}\definecolor{methcol}{RGB}{52,92,148}
\definecolor{excol}{RGB}{0,135,90}\definecolor{attncol}{RGB}{200,40,55}
\tcbset{boxbase/.style={enhanced,breakable,boxrule=0.9pt,arc=2.5pt,
  left=3mm,right=3mm,top=2.3mm,bottom=2.3mm,fonttitle=\bfseries,coltitle=white}}
% --- boîtes TUTORIEL (noms EXACTS, ne pas renommer) ---
\newtcolorbox{cle}[1][]{boxbase,colback=defcol!6,colframe=defcol,
  title={\textsc{À retenir}\ifx\relax#1\relax\relax\else\textmd{ : \itshape #1}\fi}}
\newtcolorbox{methode}[1][]{boxbase,colback=methcol!7,colframe=methcol,sharp corners,
  title={\textsc{Méthode}\ifx\relax#1\relax\relax\else\textmd{ --- \itshape #1}\fi}}
\newtcolorbox{exemplebox}[1][]{boxbase,colback=excol!5,colframe=excol,
  title={\textsc{Exemple}\ifx\relax#1\relax\relax\else\textmd{ : \itshape #1}\fi}}
\newtcolorbox{piege}[1][]{boxbase,colback=attncol!6,colframe=attncol,
  title={\textsc{Piège}\ifx\relax#1\relax\relax\else\textmd{ : \itshape #1}\fi}}
\newtcolorbox{remarque}[1][]{boxbase,colback=black!4,colframe=black!45,
  title={\textsc{Remarque}\ifx\relax#1\relax\relax\else\textmd{ : \itshape #1}\fi}}
% --- boîtes EXERCICE / CORRIGÉ (noms EXACTS, auto-comptées) ---
\newtcolorbox[auto counter]{exo}[1][]{boxbase,colback=primary!4,colframe=primary,
  title={\textsc{Exercice}~\thetcbcounter\ifx\relax#1\relax\relax\else\textmd{ --- \itshape #1}\fi}}
\newtcolorbox[auto counter]{corrige}[1][]{boxbase,colback=excol!4,colframe=excol,
  title={\textsc{Corrigé}~\thetcbcounter\ifx\relax#1\relax\relax\else\textmd{ --- \itshape #1}\fi}}
\newcommand{\R}{\mathbb{R}}\newcommand{\N}{\mathbb{N}}\newcommand{\Z}{\mathbb{Z}}\newcommand{\ds}{\displaystyle}
% (un \usepackage{fancyhdr}+en-tête est possible ; il est ignoré côté web)
% =========================================================================
\begin{document}
\newcommand{\AX}[2]{\ensuremath{\mathrm{AX_{#1}E_{#2}}}}
\begin{exo}[la résonance de \ce{ClO4-}]
Dans l'ion perchlorate, la « meilleure » structure place une charge $-1$ sur \emph{un} oxygène et
trois doubles liaisons sur les autres — mais rien ne distingue physiquement les quatre oxygènes.
\begin{enumerate}[label=\alph*)]
  \item Combien de formes limites équivalentes peut-on écrire ?
  \item Comment la charge $-1$ est-elle réellement répartie sur les oxygènes ?
  \item Que dire des quatre liaisons \ce{Cl-O} réelles ? La géométrie tétraédrique est-elle affectée ?
\end{enumerate}
\end{exo}

\begin{exo}[le benzène \ce{C6H6}]
Le benzène est un cycle de six carbones. On l'écrit avec trois doubles liaisons alternées.
\begin{enumerate}[label=\alph*)]
  \item Combien de formes limites (dites « de Kekulé ») équivalentes existe-t-il ?
  \item Que peut-on dire des six liaisons \ce{C-C} dans la molécule réelle ?
  \item En une phrase, à quoi correspond physiquement la délocalisation des électrons $\pi$ ici ?
\end{enumerate}
\end{exo}

\begin{exo}[des formes limites qui ne se valent pas]
Pour une certaine molécule, on écrit deux formes limites : dans la forme \textbf{(I)}, tous les
atomes respectent l'octet et toutes les charges formelles sont nulles ; dans la forme \textbf{(II)},
apparaît une séparation de charges $+1 / -1$.
\begin{enumerate}[label=\alph*)]
  \item Ces deux formes contribuent-elles \emph{également} à l'hybride ?
  \item Laquelle « pèse » le plus, et selon quel critère ?
\end{enumerate}
\end{exo}

\begin{exo}[corriger une idée fausse]
Un étudiant lit \ce{O=O-O <-> O-O=O} et en conclut que l'ozone est un mélange qui « bascule » sans
cesse entre deux molécules différentes. En \textbf{deux points précis}, explique pourquoi cette
interprétation est fausse.
\end{exo}

\begin{exo}[synthèse : l'ion formiate \ce{HCO2-}]
L'ion formiate a pour atome central le carbone, lié à un \ce{H} et à deux oxygènes.
\begin{enumerate}[label=\alph*)]
  \item Calcule $n_v$ et vérifie qu'on a $3$ liaisons $\sigma$ $+$ $1$ double liaison autour du carbone.
  \item Écris les deux formes limites et indique la répartition de la charge $-1$.
  \item Donne la notation \AX{n}{m} du carbone et sa géométrie. La résonance la modifie-t-elle ?
\end{enumerate}
\end{exo}

\end{document}
